\section{Discussion and Conclusion}
We presented \method, a compiler-guided agentic system for optimizing AD kernels. Given target kernels, it jointly synthesizes and optimizes forward computation, backward computation, and their saved-state interface. Users also supply a model and workload configurations for model-level verification and evaluation; \method returns selected kernel implementations and their deployment bindings. Its seed synthesizer exposes compiler-derived training semantics as editable implementations; its evolver explores the joint kernel program space; and its verifier and evaluator connect local optimization to numerical acceptance and training-level deployment selection.

Across five workflows, only 107 of 170 submissions are accepted, and only 17 of 28 accepted implementations with deployed-operator speedups retain statistically supported training gains. \method produces qualifying combined deployments on five of six evaluated workloads, but its end-to-end gains remain modest. These results expose a persistent gap in agentic training optimization: local correctness and speed do not reliably translate into correct, faster training.

Saved-state search is complementary to graph-level activation checkpointing and rematerialization: \method changes the custom forward--backward program and its internal interface, while the surrounding compiler may optimize compiler-visible values. Our results do not isolate the contribution of seed synthesis, establish long-term training stability, or characterize this interaction with subsequent compilation. Designing agentic workflows that directly optimize end-to-end training remains an opportunity for future work.
% \method demonstrates that automatic training-kernel generation is not merely
% a problem of translating a backward formula into Triton. The system must also
% construct an independent correctness contract, choose saved forward state,
% search reduction and fusion strategies, and specialize deployment across
% shape regimes. Our current benchmark schema primarily targets two-dimensional
% row-by-column operators; attention, RoPE, and fused-linear families require a
% more general multi-input and multi-output representation. In addition, task
% self-checking validates contract wiring but does not fully validate numerical
% tolerance calibration, and evolutionary runs still exhibit nontrivial
% variance. Addressing these limitations is necessary for broader automation,
% but the current results show that correctness-constrained evolution with simple
% shape dispatch can already produce robust, high-performance autograd kernels.

% \ahan{Putting down a thought here. Though I like the current challenge. In my opinion another key challenge is the interplay of correctness and optimization. The crux being, EvoGrad must output fundamentally correct kernels by construction (at the end of its optimization "passes"). However, this is challenging to do, because there are no correctness guarantees on the output of LLM code-generation. So we need to design an LLM-in-the-loop strategy that iteratively optimizes kernels in a manner that outputs a correctly code-generated kernel in the end. Another thing about this correctness uniqueness to AD is that the user codes and thinks about and codes the forward pass, whilst the backwards pass is hidden and automatically generated. Thus viable solutions must be generated for llm code-optimization. This uniqueness isn't there in a lot of other LLMs for code-optimization literature in the way its present for us (unless I am incorrect). The added layer of indirection makes this evermore challenging as the seed for LLM input must be code-generated and correct first, prior to optimization. Our design space must necessarily then consist of formal techniques + LLM optimization. This makes our design explicit: a separation of concerns between lowering and optimization, lowering is formal whilst optimization is aided by LLM, giving a ground truth to repeatedly check against. Without making correctness a first class citizen that we explicitly curated a design for, it may be tough to motivate the way we do things in EvoGrad. Just laying down some thoughts. What do folks think?} 
% \haochen{Most works I saw emphasize correctness during the optimization phase, either rejecting the failed kernels (KernelEvolve) or feed it to  a debug agent (EGG, ICML 2026) per iteration, which answers the first part of the challenge. Right now since we're using OpenEvolve, if an optimized kernel is problematic we just drop it. I think the unique challenge posed by AD which triggers our lowering + optimization design could be something new wrt correctness. }
% \minjia{This is a very good point. My thinking of this: (i) how do we evaluate the correctness and performance? If we can't even measure them, it is very hard for LLM to do this too. This goes back to the benchmark side. I believe in KernelBench, they considered this and define metrics fast-0 (correct) and fast-1 (correct and faster) for their included testing cases. We can probably learn from it to encapsulate these two properties either in the evaluation of our metric and then in our optimization pipeline. }